\chapter{Commercial identification, counterfactuals and the information a bank must supply}
\label{exec13:commercial-identification}

The preceding chapter made acquisition conditional on preserving the scientific
question. This chapter develops that condition through identification,
counterfactuals, heterogeneity, instruments, event histories and policy value.
It then compares the engine and application boundaries that carry those objects.
The method studies establish the targets and counterexamples; the later
product studies use them to choose a particular interface, entitlement and
qualification path. These are related decisions, not independent demonstrations
of the same claim.

A lending manager asks whether offering a temporary rate reduction will improve
twelve-month repayment among existing borrowers. A second manager asks whether
one particular borrower would have repaid without the offer. A third asks how
much of the effect operates through an intermediate change in payment behavior.
These requests share business language and variables, but they demand different
causal information. Purchasing a package with a causal menu does not make them
the same statistical problem. Nor does implementing a common API remove their
different assumptions.

The proposed product should first identify the requested quantity and then
compare the information and computation required to answer it. Commercial
software provides substantive precedents: graphical intervention workflows,
treatment-effect estimators, mediation models and, in at least one inspected
guide, an explicit abduction workflow. Those are capabilities against which the
proposal must be assessed. Their documentation also illustrates why a numerical
answer needs a stated model and why an unsupported question must remain visible.

This comparison draws on a declared set of sixteen commercial families and
27 bounded public document requests, together with the preserved commercial
corpus. It inspects technical documents for BayesiaLab, Stata and Bayes Server;
an Mplus resource index supplies additional leads. Other families retain
specific unresolved comparisons. The literature cutoff remains 7 October 2026.
New document copies were acquired on 13 October, and mutable pages do not
authenticate their exact contents at the cutoff. Stata's inspected PDFs carry
2025 copyright and release-19 data references; the Bayes Server guide names
specific version thresholds. These identities are more informative than an
undated search result, but they do not eliminate the historical-copy limitation.
No licensed vendor engine was executed for this comparison
\citep{exec13_ci_scope}.

\section{Four questions a commercial tool may answer}

Let $A\in\{0,1\}$ denote the offer and $Y(a)$ repayment by a specified date
under offer state $a$. The outcome is defined for every existing loan, including
loans receiving no offer. The first target is
\[
 \tau=\mathbb E\{Y(1)-Y(0)\}.
\]
The population, date, treatment version and accounting meaning are part of
this definition. A target among recipients, conditional on pre-offer history,
or under a sequence of actions needs a different expression. An analysis of
only newly approved loans would additionally face selection and possibly a
different outcome definition; it cannot silently reuse this target.

\subsection{Adjustment and general identification}

An adjustment method starts from a justified observed set $X$. Under consistency,
conditional exchangeability and suitable positivity, it evaluates
\[
 \mathbb E\{Y(a)\}
 =\int \mathbb E[Y\mid A=a,X=x]\,dP(x).
\]
The integral is an identification statement; fitting its regression and
averaging predictions is a subsequent estimation problem. A tool can implement
the latter while relying entirely on the analyst for the former.

A broader identification request supplies a graph class, observed variables,
available observational or experimental distributions and an intervention query.
For example, it may ask for $P(Y\mid\operatorname{do}(A=a))$ in a graph with
unobserved common causes, or for a conditional intervention distribution
$P(Y\mid\operatorname{do}(A=a),Z=z)$. Claims about ID or IDC must name the
algorithm and its supported graph and query classes. Support for one valid
adjustment set does not establish those broader capabilities.

The response contract should distinguish an identified functional, a justified
nonidentification result within the declared model class, and an unknown result
because the available method does not cover the request. Failure of a
back-door search belongs to the third category unless a separate argument
establishes the second. The exact mediator example below shows the practical
importance of this distinction.

\subsection{Individual counterfactuals and mediated effects}

For a borrower who received the offer and repaid, a different query is
\[
 q=\Pr\{Y(0)=0\mid A=1,Y=1\}.
\]
It concerns a relationship between the factual and alternative worlds for
the same unit. Knowing both marginal distributions $P\{Y(0)\}$ and
$P\{Y(1)\}$ need not identify their joint distribution. A structural model
with specified exogenous variables can supply that relationship, but the
relationship is then additional model information. A workflow that infers
those exogenous variables from evidence, changes an action and predicts the
result must preserve the assumptions connecting the worlds.

Mediation makes another distinction. Let $M(a)$ be a mediator, measured after
the offer, and $Y(a,m)$ the outcome under joint interventions. One decomposition
of the total effect uses
\begin{align*}
 \mathrm{NDE}&=\mathbb E[Y(1,M(0))-Y(0,M(0))],\\
 \mathrm{NIE}&=\mathbb E[Y(1,M(1))-Y(1,M(0))].
\end{align*}
Their sum is the total effect under the corresponding composition assumptions.
These are population averages of nested potential outcomes. They are neither
the same target as an individual borrower's $q$ nor automatically identified
by randomizing $A$. The mediator--outcome relation and treatment-induced
confounding still matter. Stata's one-mediator manual makes these assumptions
explicit rather than presenting the word ``counterfactual'' as an unrestricted
capability \citep{exec13_ci_stata_mediate}.

\subsection{Partial identification and uncertainty}

When the information identifies a set rather than a point, the product should
return that set and explain which restrictions would narrow it. This differs
from a confidence interval for a point parameter. A finite-sample confidence
region for an identified set requires an additional sampling argument. A
software comparison must therefore record separately whether a package computes
a point functional, logical bounds, a sensitivity region or statistical
uncertainty, and how failed or empty requests are represented.

\section{BayesiaLab: a documented graphical adjustment workflow}

The inspected graphical-identification chapter explains causal and noncausal
paths, adjustment in a supplied DAG and a particular hidden-common-cause
obstruction. Its discussion supports a comparison of adjustment semantics.
It does not specify a complete general ID/IDC implementation
\citep{exec13_ci_bayesia_criteria}. In particular, its example in which a new
hidden cause defeats adjustment should not be generalized to every graph with
hidden confounding.

The graph-surgery walkthrough is operationally more specific. An analyst draws
the causal graph, associates data columns and state values with its nodes,
estimates the probability tables, and inspects the paths to a selected target.
The intervention menu removes the treatment's incoming arcs for the causal
query. The displayed probabilities then change without treating the action as
ordinary observed evidence. This is a documented supplied-model intervention
workflow, including a visual explanation of the observational/interventional
distinction \citep{exec13_ci_bayesia_surgery}.

For the proposed bank tool, data-to-node association is a substantive capability.
An offer flag recorded after acceptance may not represent the assigned offer;
a repayment field may be available after the declared decision date. A graph
display that preserves node names but loses this mapping can give a convincing
diagram of the wrong question. The proposal should make each node's unit,
clock, state coding and source field reviewable, and make the intervention
status visible alongside the graph. These are concrete requirements suggested
by the documented commercial workflow, not evidence that the local prototype
already implements a comparable visual environment.

The likelihood-matching walkthrough documents another route for the same
example. It uses probabilistic evidence and a \emph{Fix Probabilities} operation
to hold a confounder's marginal distribution fixed while varying treatment
evidence. Its reported example matches the graph-surgery calculation. The
important mathematical question is whether the variables and distributions
being held fixed supply a valid adjustment for the requested target
\citep{exec13_ci_bayesia_matching}. A numerical balancing operation does not
establish that an unmeasured common cause has been removed.

Bayesia's contribution-analysis page uses generated data, a learned network,
expert assumptions about confounder selection and Likelihood Matching to compare
factual and intervention conditions. That supports a documented contribution
workflow. The inspected text does not fully specify a general individual-level
abduction model or the coupling of alternative outcomes. Its linked slides
and recording were not inspected in this comparison
\citep{exec13_ci_bayesia_counterfactual}. The proposed product should therefore
ask whether a request is a population intervention contrast, a contribution
allocation or a conditional individual counterfactual before choosing an adapter.

The reuse recommendation is conditional: treat these workflows as commercial
benchmarks and possible licensed integration targets. Obtain the exact API,
edition, export semantics and deployment rights before promising an adapter.
Porting proprietary implementation is not an available reuse strategy on the
basis of public documentation. Independently implementing a justified
mathematical method requires its own provenance and validation.

\section{Stata: treatment effects and mediation with explicit assumptions}

The inspected \texttt{teffects} introduction documents average treatment effects,
effects among treated subjects and potential-outcome means. It covers regression
adjustment, inverse-probability weighting, augmented weighting, weighted
regression adjustment and matching. Its caveats discuss independent sampling,
conditional independence and overlap, including why unmeasured determinants of
treatment and outcome can invalidate an estimate
\citep{exec13_ci_stata_causal_intro}.

This is a substantial estimator baseline. The proposal needs a defensible reason
to build another estimator rather than a thin, licensed workflow around an
existing one. A possible reason is preservation of the business target,
information clock, graph ambiguity and source provenance across several tools.
That reason becomes persuasive only when a minimal prototype demonstrates the
behavior and a comparison measures the total effort of obtaining an acceptable
answer. Merely presenting additional estimators or claiming a common CLI does
not establish superiority over an established analysis workflow.

For the lending request, an analyst could supply pre-offer covariates and fit
one of the documented treatment-effect estimators. The output then inherits
the asserted covariate sufficiency and the relevant sampling law. The selected
manual's IID discussion does not establish validity for every repeated monthly
loan record or shared-calendar shock. Likewise, an overlap diagnostic concerns
observed treatment support; it cannot certify that an omitted repayment driver
has been measured. The proposed wrapper must record these distinct obligations
instead of converting a successful command exit into an identification pass.

The \texttt{mediate} manual supplies a separate, developed capability. It
describes continuous, binary and count outcome/mediator models, several treatment
types, interactions, direct and indirect effects, potential-outcome means and
variance options. It explicitly addresses unmeasured confounding in the
treatment--outcome, treatment--mediator and mediator--outcome relations, and
excludes treatment-induced mediator--outcome confounding under its stated
identification conditions \citep{exec13_ci_stata_mediate}.

Consider an offer that changes a borrower's monthly payment burden, which then
affects repayment. The offer might also trigger a collections contact that
affects both payment behavior and repayment. That contact is a possible
treatment-induced mediator--outcome confounder. Adding it mechanically to every
regression is not a justification of the natural-effect estimand. The business
request may need a different mediated target or a richer design. A common API
should carry the precise nested expression and its assumptions rather than
label every coefficient product an indirect causal effect.

The inspected manuals support these estimator and mediation comparisons.
They do not establish generic graphical ID/IDC, arbitrary individual
abduction or complete latent-graph coverage. Other Stata commands may address
other endogeneity models; those require their own technical comparisons. The
current product proposal should prioritize matched documented targets and keep
unreviewed extensions explicit. Exact licensing, embedded execution and price
comparisons remain outstanding; no commercial runtime result is reported here.

\section{Bayes Server: intervention, abduction and the need for an exact API contract}

The Bayes Server causal guide distinguishes ordinary evidence from intervention,
describes discrete and continuous variables, and gives a concrete individual
counterfactual workflow. Its example enters factual evidence, uses an
\emph{Abduction} dialog to infer characteristic variables, applies a changed
intervention, and reads the resulting outcome. The guide describes helper
functions from version 10 onward and says the example calculations can be
performed from version 9.4 \citep{exec13_ci_bayesserver_causality}.

This is stronger evidence of a documented abduction workflow than a generic
``what if'' feature label. It reveals another capability gap for the proposal:
the local reference service does not yet offer this general structural-model
workflow. The architecture should represent the exogenous-variable model,
conditioning evidence, action and resulting distribution as separate objects.
It should also retain whether abduction produced a unique value, a distribution
or an unresolved set. A model-specified answer can be useful while remaining
conditional on the model's cross-world assumptions.

The guide also lists back-door and front-door adjustment, instruments,
disjunctive-cause reasoning and do-calculus, stating that many approaches are
available from version 10 onward. That phrasing does not identify which exact
algorithms are implemented, the graph classes for which they are complete, or
the supported failure certificates. Section~\ref{exec14:commercial-contracts} develops the retained Advanced
Causal Query and Abduction guides and API signatures. Their criterion-specific
information resolves this narrower documentation question; complete ID/IDC
and arbitrary counterfactual identification remain unestablished.

An adapter proposal needs more than the guide's terminology. It needs exact
node and edge encodings, latent-variable treatment, allowed interventions,
conditioning behavior, numerical types, export format, method/version identity
and a licensed execution boundary. Until those are inspected and tested, the
recommendation is a focused commercial benchmark candidate. It is not a promise
of immediately reusable functionality or a claim that a general causal
identification engine is absent from the product.

\section{The remaining commercial landscape and explicit omissions}

The Mplus mediation index names counterfactual causal effects, examples and
\texttt{MODEL INDIRECT} material associated with version 7.2. It is a useful
route to original technical sources \citep{exec13_ci_mplus_mediation}. An index
does not supply the assumptions or proofs of the linked papers. Those originals
must be acquired as genuine PDFs and inspected before they support additional
scholarly claims. The existing latent-variable discussion also does not settle
this distinct identification question.

The table records the original acquisition boundaries. ``Unresolved'' describes
the evidence available to this survey. It does not mean the vendor lacks the
capability. Price and edition entitlements have not been verified for the
proposed embedding scenario \citep{exec13_ci_scope}.

\begin{longtable}{P{0.17\textwidth}P{0.34\textwidth}P{0.39\textwidth}}
\toprule
Family & Available comparison & Specific remaining source or decision \\
\midrule
\endhead
BayesFusion & Preserved structure, pattern and license documents; new causal route unresolved. & Exact intervention/query semantics and wrapper/edition support. \\
HUGIN & The initial request failed; the later GUI 9.8 manual supplies decision/information and model-class boundaries. & General identification/query semantics and the selected API/edition remain to be established. \\
Netica & Preserved learning documentation concerns adjacent operations. & Intervention and counterfactual semantics, method boundaries and license. \\
SAS & Treatment/mediation PDF transfers were incomplete; replacement HTML routes failed. & Authentic CAUSALTRT/CAUSALMED technical text and target-specific assumptions. \\
IBM & Preserved temporal-model material does not answer general identification. & Edition-specific causal/Amos methods; distinguish SEM fitting from causal identification. \\
MathWorks & Initial access was denied; later R2026b panel-method pages provide an adjacent estimation comparison. & Treatment/identification documentation with target, assumptions and runtime scope remains distinct. \\
causaLens & The documentation endpoint redirected to login. & Public technical material or subsequently authorized copies; no inference from the login page. \\
EViews & Earlier structural-model navigation supplies limited adjacent context. & Actual restrictions, intervention/scenario methods and API contract. \\
OxMetrics & PcGive/PcGets/Autometrics leads remain incompletely resolved. & Versioned identification, intervention and uncertainty methods. \\
LISREL & Retained as a structural-model candidate. & Primary model restrictions, identification diagnostics and embedding rights. \\
JMP & Earlier routes failed; retained treatment documentation is developed below. & Preserve target and edition; installed automation and rights remain unresolved. \\
Analytica & Two attempted official routes returned missing-page responses. & Influence-diagram and Decision Engine semantics, separating supplied decisions from identified effects. \\
\bottomrule
\end{longtable}

The following studies resolve parts of this initial inventory: JMP receives
a developed treatment comparison, and GeNIe/SMILE receives a supplied-model
decision and engine comparison. Those findings do not establish unrestricted
graphical identification. Chapter~\ref{scisyn:software} separately records a
named causaLens service offering; its advertised scope does not resolve the
versioned identification contract sought here. Section~\ref{gap29:software-followup}
also develops the later HUGIN decision-model and MATLAB panel-method evidence.
HUGIN's information links and discrete chance-node restriction matter to a
servicing comparison; MATLAB's panel transformations and subject covariance
matter to a branch study. Neither finding establishes a general identification
engine. Their retrieval after the literature cutoff is retained explicitly.
\citep{gap29hugin,gap29matlab}
The remaining omissions make
complete commercial coverage premature, while allowing a useful proposal
recommendation. The immediate recommendation can
use the three developed comparisons, give concrete acquisition actions for
the others, and select public exact fixtures for the next prototype. A bank
experiment is unnecessary merely to establish what a manual says. A paid trial
is also unnecessary for the mathematical reference examples below.

\section{Exact lending examples that separate the required capabilities}

The following are project constructions, evaluated with exact rational
arithmetic. They are neither fitted vendor outputs nor bank observations.
Their purpose is to make identification, partial identification and
counterfactual semantics reviewable before an expensive experiment
\citep{exec13_ci_examples}.

\subsection{The same observed loan table can support different effects}

Consider an observed law
\[
\begin{array}{c|cc}
 &Y=0&Y=1\\\hline
A=0&3/8&1/8\\
A=1&1/8&3/8
\end{array}
\]
with all four cells positive. The observed repayment difference is
$3/4-1/4=1/2$. One generating model has a binary pre-existing repayment type
$U$ with probability $1/2$, an independent bit $V$ with probability $1/4$,
offer $A=U\mathbin{\mathrm{xor}}V$, and outcome $Y=U$. Changing the offer
does not change repayment in this model: $\tau=0$.

A second model randomizes $A$ with probability $1/2$. With an independent
uniform four-state variable $W\in\{0,1,2,3\}$, let
$Y(0)=\mathbf 1\{W<1\}$ and $Y(1)=\mathbf 1\{W<3\}$.
Its effect is $\tau=1/2$, yet its observed table is identical. Enumerating
both finite models verifies the equality exactly. Additional rows from the
same observed law cannot choose between them. Observed positivity does not
resolve hidden confounding.

If the analyst commits to the second model's randomized assignment, a
treatment-effect estimate has a causal interpretation. If the information
allows both models, the same calculation is an association with unresolved
causal meaning. A commercial adapter and the local service should preserve
that difference even when they display the same number.

\subsection{Sharp bounds remain useful information}

With only consistency and binary potential outcomes, the observed successful
treated borrowers contribute $3/8$ to $\mathbb E[Y(1)]$. The unobserved
$Y(1)$ values among the untreated can contribute anything from zero to $1/2$.
Similarly, observed successful untreated borrowers contribute $1/8$ to
$\mathbb E[Y(0)]$. Thus
\[
 \mathbb E[Y(1)]\in[3/8,7/8],\qquad
 \mathbb E[Y(0)]\in[1/8,5/8],\qquad
 \tau\in[-1/4,3/4].
\]
These bounds are jointly attainable. Allocate each of the four observed
$(A,Y)$ cells to either value of its missing potential outcome. The $2^4=16$
resulting vertices preserve the table and attain both effect endpoints.
Every fractional allocation is a convex combination of the vertices, so a
linear effect cannot exceed these extrema. The exact program checks this
enumeration separately from the missing-outcome algebra.

For a hypothetical planning value of USD~1,000 per incremental repayment and
a USD~20 cost per offer, net value is $1000\tau-20$. The two equal-observed-law
models imply respectively minus USD~20 and plus USD~480 per offer. The
consistency-only range is minus USD~270 to plus USD~730. This is an assumed
value mapping, not a measured ROI or a full loan-accounting model. Its purpose
is to show why the unresolved scientific distinction can reverse a decision.

The range includes both profitable and unprofitable offers. It can justify
considering a better information design, but it does not by itself establish
that any particular experiment is worth its cost. That further argument needs
the experiment's expected resolution, the affected decisions and their costs.

\subsection{Hidden confounding does not always prevent identification}

Now assume the precise structural graph $U\to A$, $U\to Y$,
$A\to M\to Y$, with no additional arrows and independent disturbances.
Let
\begin{align*}
 \Pr(U=1)&=1/2,\\
 \Pr(A=1\mid U=u)&=1/4+u/2,\\
 \Pr(M=1\mid A=a)&=1/5+3a/5,\\
 \Pr(Y=1\mid M=m,U=u)&=1/10+2m/5+2u/5.
\end{align*}
The mediator intercepts all directed effects of $A$ on $Y$ in this stipulated
law. The observed functional
\[
 \mu_a=\sum_m P(m\mid a)\sum_b P(Y=1\mid m,b)P(b)
\]
equals the intervention mean here. To verify rather than assume this equality,
note that $U$ and $M$ are conditionally independent given $A$, so
\[
 P(Y=1\mid m,b)=1/5+2m/5+b/5.
\]
Averaging over $P(A=b)=1/2$ gives $3/10+2m/5$; averaging over
$P(M=1\mid A=a)$ gives $19/50+6a/25$.
Independently summing the intervened structural law over $(U,M)$ produces
the same two means, $19/50$ and $31/50$. The effect is $6/25=0.24$,
while the ordinary observed difference is $11/25=0.44$.

The example is a useful front-door fixture, not a guarantee for every mediator.
It would require another argument if the offer had a direct effect on repayment,
if an unmeasured cause affected both offer and mediator, or if the measured
mediator failed to represent the stipulated variable. The immediate software
requirement is nonetheless clear: a failed observed back-door adjustment must
not be reported as a universal nonidentification proof. Under the hypothetical
USD~1,000/20 mapping, this model's identified net value is USD~220 per offer;
the model assumptions carry that interpretation.

\subsection{Even complete randomized margins can leave an individual question open}

Let $A$ and $U$ be independent fair bits. In one world set
$Y(0)=Y(1)=U$; in a second set $Y(0)=U$ and $Y(1)=1-U$.
Both worlds have $P\{Y(a)=1\}=1/2$ for each action, the same randomized
observed law, and average effect zero. In the first world, a treated borrower
who repaid would also have repaid without the offer, so $q=0$. In the second,
that borrower would not have repaid, so $q=1$.

The same margins allow the population probability of benefit
$P\{Y(1)=1,Y(0)=0\}$ to range from zero to $1/2$. A randomization result
can identify an average effect while leaving this cross-world association
unresolved. An abduction engine supplied with one complete structural model
may return an answer, but the evidence must say which coupling it used. A
language model's plausible borrower narrative is not additional identifying
information.

\section{A proposal architecture that can be compared fairly}

The smallest coherent prototype should expose a query object containing the
population, outcome clock, intervention, conditioning event and target units.
It should attach the supplied graph or structural model, observed-variable
map, permitted data laws and assumptions. The method registry should state
the exact domain of each method. A response should then carry its expression
or set, derivation/certificate where available, unresolved premises and the
next permissible operation.

Estimation belongs after an identified functional has been selected. A
licensed Stata adapter, a documented graphical intervention adapter, or an
open-source estimator could occupy that stage. The comparison should hold the
query, data, adjustment set and sampling assumptions fixed. Method-specific
diagnostics, missing values and failed requests remain part of the result.
Numerical agreement is meaningful only when both tools answered the same
question under the same assumptions.

For individual counterfactuals, a separate structural-model path should retain
the posterior exogenous distribution after abduction and the exact action used
in prediction. It should refuse to imply unique individual identification when
the available observations permit several couplings. For partial identification,
the interface should preserve closed/open endpoints, unions and empty sets
where the method requires them. These are proposed requirements; the local
back-door reference implementation and deterministic workflow checks establish
only their documented subset.

The exact fixtures above permit a useful next experiment with no authentic
bank data. First compare each proposed method's supported-domain response to
the independently derived truth: equal observed laws, a front-door functional,
sharp bounds and differing cross-world couplings. Then add finite public or
generated samples under a frozen observation law, with uncertainty, availability
and width criteria declared before final draws. A vendor runtime comparison
can follow once technical contracts and licensing are resolved. Neither an
unavailable commercial manual nor a deferred paid trial needs to block the
public prototype route.

The investment case should consequently distinguish three decisions: adopt
an existing estimator/workflow where its domain and rights match; build the
missing scientific contract and cross-tool explanation where that provides
demonstrable value; and defer claims requiring undocumented or unidentified
capabilities. This makes the commercial landscape consequential for the
proposal while retaining honest limits on source coverage and implementation.
The next study makes the identification and abduction contracts more concrete;
the subsequent estimator and product comparisons retain the same distinction
between a defined target, an implemented operation and justified uncertainty.
\label{exec13:commercial-identification-end}
